% ==================== 宏包 ====================
\usepackage[top=2.54cm,bottom=2.54cm,left=3.17cm,right=3.17cm]{geometry}
\usepackage{setspace}
\setstretch{1.5}
\usepackage{amsmath,amssymb,amsfonts}
\usepackage{amsthm}
\let\emptyset\varnothing
\usepackage{graphicx}
\usepackage{adjustbox}
\usepackage{booktabs}
\usepackage{tabularx}
\usepackage{multirow}
\usepackage{longtable}
\usepackage{enumitem}
\setlist[enumerate,1]{label=（\arabic*）,wide=2em,labelsep=0pt,listparindent=2em,parsep=0pt,topsep=0pt,partopsep=0pt}
\usepackage{algorithm2e}
\usepackage{listings}
\usepackage{xcolor}
\usepackage[labelsep=space]{caption}
\captionsetup{font={normalsize,bf}, labelsep=space, aboveskip=0pt, belowskip=0pt}
\usepackage{subcaption}
\usepackage{float}
\usepackage{url}
\usepackage{tikz}
\usetikzlibrary{shapes.geometric, arrows.meta, positioning, fit, backgrounds, calc, patterns, decorations.pathreplacing}
\usepackage{pgfplots}
\pgfplotsset{compat=1.18}
\usepgfplotslibrary{fillbetween, polar, patchplots, colormaps}
% 实验图 csv 数据搜索路径（论文 cwd 为 毕业论文/，fig 内引用形式 data/<scn>/xxx.csv）
\pgfplotsset{table/search path={../图片/, ../, .}}
\usepackage[backend=biber,style=gb7714-2015,gbpub=false,
            uniquename=init,uniquelist=true,
            maxcitenames=2,mincitenames=1,
            sorting=none,doi=false,url=false]{biblatex}
\usepackage{titletoc}

% ==================== 字体设置 ====================
\def\ProjectFontDir{../fonts}
% Project-local font setup for XeLaTeX.
% Callers may set \ProjectFontDir before \input'ing this file.
\providecommand{\ProjectFontDir}{../fonts}

\setmainfont[
  Path = \ProjectFontDir/tex-gyre/,
  Extension = .otf,
  UprightFont = *-regular,
  BoldFont = *-bold,
  ItalicFont = *-italic,
  BoldItalicFont = *-bolditalic,
  Ligatures = TeX
]{texgyretermes}

\setCJKmainfont[
  Path = \ProjectFontDir/fandol/,
  Extension = .otf,
  BoldFont = FandolSong-Bold,
  ItalicFont = FandolKai-Regular
]{FandolSong-Regular}
\setCJKsansfont[
  Path = \ProjectFontDir/fandol/,
  Extension = .otf,
  BoldFont = FandolHei-Bold,
  ItalicFont = FandolKai-Regular
]{FandolHei-Regular}
\setCJKmonofont[
  Path = \ProjectFontDir/fandol/,
  Extension = .otf,
  ItalicFont = FandolKai-Regular
]{FandolFang-Regular}

\setCJKfamilyfont{zhsong}[
  Path = \ProjectFontDir/fandol/,
  Extension = .otf,
  BoldFont = FandolSong-Bold
]{FandolSong-Regular}
\setCJKfamilyfont{zhsongb}[
  Path = \ProjectFontDir/fandol/,
  Extension = .otf
]{FandolSong-Bold}
\setCJKfamilyfont{zhhei}[
  Path = \ProjectFontDir/fandol/,
  Extension = .otf,
  BoldFont = FandolHei-Bold,
  ItalicFont = FandolKai-Regular
]{FandolHei-Regular}
\setCJKfamilyfont{zhheib}[
  Path = \ProjectFontDir/fandol/,
  Extension = .otf
]{FandolHei-Bold}
\setCJKfamilyfont{zhfs}[
  Path = \ProjectFontDir/fandol/,
  Extension = .otf
]{FandolFang-Regular}
\setCJKfamilyfont{zhkai}[
  Path = \ProjectFontDir/fandol/,
  Extension = .otf
]{FandolKai-Regular}

\ExplSyntaxOn
\cs_if_exist:NT \ctex_punct_set:n
  {
    \ctex_punct_set:n { fandol }
    \ctex_punct_map_family:nn   { \CJKrmdefault         } { zhsong  }
    \ctex_punct_map_family:nn   { \CJKsfdefault         } { zhhei   }
    \ctex_punct_map_family:nn   { \CJKttdefault         } { zhfs    }
    \ctex_punct_map_bfseries:nn { \CJKrmdefault, zhsong } { zhsongb }
    \ctex_punct_map_bfseries:nn { \CJKsfdefault, zhhei  } { zhheib  }
    \ctex_punct_map_itshape:nn  { \CJKrmdefault         } { zhkai   }
  }
\ExplSyntaxOff

\providecommand{\songti}{}
\providecommand{\heiti}{}
\providecommand{\fangsong}{}
\providecommand{\kaishu}{}
\renewcommand{\songti}{\CJKfamily{zhsong}}
\renewcommand{\heiti}{\CJKfamily{zhhei}}
\renewcommand{\fangsong}{\CJKfamily{zhfs}}
\renewcommand{\kaishu}{\CJKfamily{zhkai}}

% 圆圈数字 ①-⑳ 等 (U+2460-U+24FF) 强制走 CJK 字体，避免 TeX Gyre Termes Bold 缺字符渲染为空
\xeCJKDeclareCharClass{CJK}{"2460->"24FF}

% ==================== 页眉页脚 ====================
\usepackage{fancyhdr}
\fancypagestyle{plain}{%
  \fancyhf{}%
  \fancyfoot[C]{\zihao{5}\rmfamily\thepage}%
  \renewcommand{\headrulewidth}{0pt}%
  \renewcommand{\footrulewidth}{0pt}%
}
\pagestyle{fancy}
\fancyhf{}
\fancyfoot[C]{\zihao{5}\rmfamily\thepage}
\renewcommand{\headrulewidth}{0pt}
\renewcommand{\footrulewidth}{0pt}

\usepackage[unicode=true,pdfencoding=unicode,psdextra=true]{hyperref}

% ==================== TikZ 统一配色 ====================
\definecolor{deepblue}{RGB}{81, 132, 178}
\definecolor{lightblue}{RGB}{170, 212, 248}
\definecolor{skyblue}{RGB}{154, 220, 255}
\definecolor{inputyellow}{RGB}{255, 248, 154}
\definecolor{outputpink}{RGB}{255, 178, 166}
\definecolor{improvedpink}{RGB}{255, 138, 174}
\definecolor{lightpink}{RGB}{241, 167, 181}
\definecolor{deeppink}{RGB}{213, 82, 118}
\definecolor{bglight}{RGB}{242, 245, 250}
\definecolor{lightgreen}{RGB}{180, 230, 180}
\definecolor{lightorange}{RGB}{255, 210, 140}
\definecolor{dangerred}{RGB}{230, 100, 100}

% 视觉增强：高饱和度对比色（实验图 MIKU vs Baseline）
\definecolor{vibrantgreen}{RGB}{46, 175, 80}
\definecolor{vibrantorange}{RGB}{240, 130, 30}
\definecolor{vibrantred}{RGB}{220, 50, 50}

% tikz 样式中用到的颜色别名（允许作为 xcolor 颜色参与混合运算）
\colorlet{obsstatic}{dangerred!70}
\colorlet{bandok}{lightgreen}
\colorlet{obsbuf}{dangerred!70}
\colorlet{egopt}{deepblue}

% 灰度打印版开关（make thesis-print 通过 jobname 激活）
\newif\ifprintversion
\ifdefined\PrintVersionFlag \printversiontrue \fi
\ifprintversion % 灰度打印版颜色覆盖 — 由 Makefile 在 \printversiontrue 时 \input
% 填充色 → 白/浅灰（保证文字可读）
\colorlet{bglight}{white}
\colorlet{lightblue}{gray!15}
\colorlet{skyblue}{gray!15}
\colorlet{inputyellow}{gray!10}
\colorlet{outputpink}{gray!20}
\colorlet{lightpink}{gray!20}
\colorlet{lightgreen}{gray!15}
\colorlet{lightorange}{gray!25}
% 强调/边框色 → 黑/深灰
\colorlet{deepblue}{black}
\colorlet{deeppink}{gray!35}
\colorlet{dangerred}{gray!55}
\colorlet{vibrantred}{gray!55}
\colorlet{improvedpink}{gray!40}
% 曲线色 — deepblue(Baseline)=黑, vibrantgreen(MIKU)=深灰, 确保可区分
\colorlet{vibrantgreen}{gray!45}
\colorlet{vibrantorange}{gray!65}
% 别名跟随
\colorlet{obsstatic}{gray!50}
\colorlet{bandok}{gray!15}
\colorlet{obsbuf}{gray!50}
\colorlet{egopt}{black}
 \fi

% 统一 tikz 样式（与 图片/figures.tex 保持同步）
\tikzset{
  mod/.style            = {rectangle, rounded corners=2pt, draw=deepblue, fill=lightblue, minimum height=0.9cm, align=center, font=\small},
  modimp/.style         = {rectangle, rounded corners=2pt, draw=deeppink, fill=improvedpink, minimum height=0.9cm, align=center, font=\small},
  modin/.style          = {rectangle, draw=deepblue, fill=inputyellow, minimum height=0.9cm, align=center, font=\small},
  modout/.style         = {rectangle, draw=deeppink, fill=outputpink, minimum height=0.9cm, align=center, font=\small},
  axx/.style            = {->, thick, black!80},
  flow/.style           = {-Stealth, thick, deepblue},
  flowimp/.style        = {-Stealth, thick, deeppink},
  obsstatic/.style      = {fill=dangerred!70, draw=dangerred, thick},
  obsdyn/.style         = {fill=lightorange, draw=dangerred, thick},
  obsbuf/.style         = {draw=dangerred!70, dashed, thick},
  bandok/.style         = {fill=lightgreen, opacity=0.55},
  bandbad/.style        = {fill=dangerred!30, opacity=0.55},
  egopt/.style          = {fill=deepblue, draw=deepblue},
  reflinestyle/.style   = {deepblue, very thick},
}

% ==================== 换行与 texttt 自动换行 ====================
\emergencystretch=3em
\tolerance=2000
\hbadness=10000
\usepackage{seqsplit}
\let\origtexttt\texttt
\renewcommand{\texttt}[1]{\origtexttt{\seqsplit{#1}}}

% ==================== 调试目录开关（图/表/代码目录，正式提交前设为 false） ====================
\newif\ifshowdebug
% \showdebugtrue  % true: 显示图/表/代码目录；false: 关闭

% ==================== 超链接配置 ====================
\hypersetup{
    colorlinks=true,
    linkcolor=black,
    citecolor=black,
    urlcolor=black,
}

\addbibresource{references.bib}

% ==================== 章节标题样式 ====================
\ctexset{
    chapter = {
        format = \raggedright\zihao{4}\heiti\bfseries,
        name = {},
        number = \arabic{chapter},
        beforeskip = 24pt,
        afterskip = 18pt,
    },
    section/format = \raggedright\zihao{-4}\heiti\bfseries,
    subsection/format = \raggedright\zihao{-4}\songti\bfseries,
    contentsname = {目\hspace{2em}录},
}

% ==================== 目录样式 ====================
\setcounter{tocdepth}{1}

\titlecontents{chapter}
  [0em]
  {\zihao{-4}\heiti\bfseries}
  {\thecontentslabel\hspace{1em}}
  {}
  {\titlerule*[0.5pc]{.}\contentspage}

\titlecontents{section}
  [2em]
  {\zihao{-4}\songti}
  {\thecontentslabel\hspace{1em}}
  {}
  {\titlerule*[0.5pc]{.}\contentspage}

% ==================== 参考文献字体 ====================
\renewcommand{\bibfont}{\zihao{-4}}

% ==================== 代码列表样式 ====================
\renewcommand{\lstlistingname}{代码}
\AtBeginDocument{\renewcommand{\thelstlisting}{\thechapter-\arabic{lstlisting}}}
\lstset{
    basicstyle=\linespread{0.85}\selectfont\zihao{-5}\ttfamily,
    keywordstyle=\color{deepblue},
    commentstyle=\color{lightgreen!60!black},
    stringstyle=\color{dangerred!70!black},
    numbers=left,
    numberstyle=\tiny\color{gray},
    frame=single,
    breaklines=true,
    columns=flexible,
    showstringspaces=false,
    language=C++,
    tabsize=4,
    aboveskip=6pt,
    belowskip=6pt,
}
% listings 不自带 json 语法，本项目 Apollo pipeline.pb.txt 配置文件按 JSON 风格高亮
\lstdefinelanguage{json}{
    keywords={true,false,null,stage,task,name,type},
    keywordstyle=\color{deepblue},
    sensitive=true,
    comment=[l]{//},
    morecomment=[s]{/*}{*/},
    commentstyle=\color{lightgreen!60!black},
    morestring=[b]",
    stringstyle=\color{dangerred!70!black},
    showstringspaces=false,
}

\BeforeBeginEnvironment{lstlisting}{\par\noindent\begin{minipage}{\linewidth}}
\AfterEndEnvironment{lstlisting}{\end{minipage}\par\medskip}

% ==================== 算法环境本地化 ====================
\RestyleAlgo{boxruled}
\SetAlgorithmName{算法}{算法}{算法列表}
\SetKwInput{KwIn}{输入}
\SetKwInput{KwOut}{输出}
\SetKw{KwRet}{返回}

\renewcommand{\algorithmautorefname}{算法}

% ==================== 定理环境 ====================
\theoremstyle{definition}
\newtheorem{lemma}{引理}
\newtheorem{theorem}{定理}
